\subsection{DIMENSION OF VECTOR SPACES}

\BourbakiRunInHeading{THEOREM 3. Two bases of the same vector space E over a field K are equipotent.}

We note first that if E admits an infinite basis B, it follows from § 1, no. 12, Corollary 2 to Proposition 23 that every other basis of E is equipotent to B. We may therefore confine our attention to the case where E has a finite basis of n elements. We note that every monogenous vector space over K, not reduced to 0, is a simple K-module (I, § 4, no. 4, Definition 7), for it is generated by each of its elements \(\neq0\) , by virtue of the relation \(\mu a = (\mu \lambda)(\lambda^{-1}a)\) for \(\mu \in K\) , \(\lambda \in K\) and \(\lambda \neq 0\) . Hence if \((a_{i})_{1 \leqslant i \leqslant n}\) is a basis of E, then \(E = \bigoplus_{i=1}^{n} Ka_{i}\) up to isomorphism and the subspaces \(E_{k} = \bigoplus_{i=1}^{k} Ka_{i}\) for \(0 \leqslant k \leqslant n\) form a Jordan-Hölder series of E, \(E_{k}/E_{k-1}\) being isomorphic to \(Ka_{k}\) . Theorem 3 then follows in this case from the Jordan-Hölder Theorem (I, § 4, no. 7, Theorem 6).

A proof can be given independent of the Jordan-Hölder Theorem, by showing by induction on \(n\) that, if \(E\) admits a basis of \(n\) elements, every other basis \(B'\) has at most \(n\) elements. The proposition is obvious for \(n = 0\) . If \(n \geqslant 1\) , \(B'\) is non-empty; then let \(a \in B'\) . By Theorem 2 (no. 1) there exists a subset \(C\) of \(B\) such that \(\{a\} \cup C\) is a basis of \(E\) and \(a \notin C\) , since \(\{a\} \cup B\) is obviously a generating system of \(E\) . As \(B\) is a basis of \(E\) , \(C = B\) is impossible (no. 1, Corollary to Theorem 2) and hence \(C\) has at most \(n - 1\) elements. Let \(V\) be the subspace generated by \(C\) and \(V'\) the subspace generated by \(B' - \{a\}\) ; \(V\) and \(V'\) are both supplementary to the subspace \(Ka\) of \(E\) and hence are isomorphic (§ 1, no. 10, Proposition 13). As \(V\) admits a basis with at most \(n - 1\) elements, \(B' - \{a\}\) has at most \(n - 1\) elements by the induction hypothesis and hence \(B'\) has at most \(n\) elements.

DEFINITION 1. The dimension of a vector space E over a field K, denoted by dim \(_{K}\) E or {[}E:K{]} (or simply dim E) is the cardinal of any of the bases of E. If M is a subset of E, the rank of M (over K), denoted by rg M or rg \(_{K}\) M, is the dimension of the vector subspace of E generated by M.

To say that E is finite-dimensional is equivalent to saying that E is a K-module of finite length and \(\dim_{K}E = long_{K}E\) .

COROLLARY. For every subset M of E, the rank of M is at most equal to dim E.

If V is the vector subspace of E generated by M, M contains a basis \(B'\) of V (no. 1, Theorem 2) and as \(B'\) is a free subset of E, it is contained in a basis B of E (no. 1, Theorem 2); then \(\text{Card}(B') \leqslant \text{Card}(B)\) , whence the corollary.

Theorems 2 and 3 immediately imply the following proposition:

PROPOSITION 1. (i) For a left vector space over \(\mathbf{K}\) to be of finite dimension \(n\) , it is necessary and sufficient that it be isomorphic to \(\mathbf{K}_s^n\) .

\begin{enumerate}
\def\labelenumi{(\roman{enumi})}
\setcounter{enumi}{1}
\item
  For two vector spaces \(\mathbf{K}_s^m\) and \(\mathbf{K}_s^n\) to be isomorphic ( \(m\) and \(n\) integers \(\geqslant 0\) ), it is necessary and sufficient that \(m = n\) .
\item
  In a vector space \(\mathbf{E}\) of finite dimension \(n\) , every generating system has at least \(n\) elements; a generating system of \(\mathbf{E}\) with \(n\) elements is a basis of \(\mathbf{E}\) .
\item
  In a vector space \(\mathbf{E}\) of finite dimension \(n\) , every free subset has at most \(n\) elements; a free subset with \(n\) elements is a basis of \(\mathbf{E}\) .
\end{enumerate}

PROPOSITION 2. Let \((\mathbf{E}_i)_{i\in I}\) be a family of vector spaces over K. Then

\[
\dim_ {\mathbb {K}} \left(\bigoplus_ {i \in I} E _ {i}\right) = \sum_ {i \in I} \dim_ {\mathbb {K}} E _ {i}.\tag{1}
\]

If the \(E_{i}\) are canonically identified with subspaces of \(E = \bigoplus_{i \in I} E_{i}\) and \(B_{i}\) is a basis of \(E_{i} (i \in I)\) , then \(B = \bigcup_{i \in I} B_{i}\) is a basis of E (§1, no. 11, Proposition 19); whence relation (1) since the \(B_{i}\) are pairwise disjoint.

Remark. (1) Examples can be given of modules admitting two finite bases not having the same number of elements (§ 1, Exercise 16(c)). However:

PROPOSITION 3. Let A be a ring such that there exists a homomorphism \(\rho\) of A into a field D; then for every free A-module E, any two bases of E are equipotent.

Consider the vector space \(\rho^{*}(\mathbf{E}) = \mathbf{D} \otimes_{\mathbf{A}} \mathbf{E}\) over \(\mathbf{D}\) obtained by extending the ring of scalars to \(\mathbf{D}\) ( \(\S 5\) , no. 1) and let \(\phi: x \mapsto 1 \otimes x\) be the canonical mapping of \(\mathbf{E}\) into \(\rho^{*}(\mathbf{E})\) ; if \((a_{\lambda})\) is a basis of \(\mathbf{E}\) , \((\phi(a_{\lambda}))\) is a basis of \(\rho^{*}(\mathbf{E})\) ( \(\S 5\) , no. 1, Proposition 4); the proposition then follows from Theorem 3.

COROLLARY. If A is a commutative ring \(\neq0\) and E a free A-module, any two bases of E are equipotent.

There exists in A at least one maximal ideal m (I, § 8, no. 6, Theorem 1) and, as A/m is a field, the conditions of Proposition 3 are fulfilled.

Remarks. (2) When a free A-module E is such that any two bases of E are equipotent, the cardinal of an arbitrary basis of E over A is also called the dimension or rank of E and denoted by \(\dim_{A}E\) or \(\dim E\) .

\begin{enumerate}
\def\labelenumi{(\arabic{enumi})}
\setcounter{enumi}{2}
\tightlist
\item
  Let A be a ring such that any two bases of a free A-module are equipotent and let K be a subfield of A, so that A can be considered as a left vector space over K by restricting the scalars. Every free A-module E can similarly be considered as a left vector space over K and it then follows from §1, no. 13, Proposition 25 that
\end{enumerate}

\[
\dim_ {\mathbb {K}} \mathrm{E} = \dim_ {\mathbb {K}} \mathrm{E}. \dim_ {\mathbb {K}} \mathrm{A} _ {s ^ {*}}\tag{2}
\]

\begin{enumerate}
\def\labelenumi{(\arabic{enumi})}
\setcounter{enumi}{3}
\tightlist
\item
  In Chapter VIII we shall see examples of rings satisfying the conclusion of Proposition 3 but not the hypothesis.
\end{enumerate}
% semantic-fragments: legacy-input-seam
\relax{}
